\documentclass{scrartcl}
\usepackage{graphicx}  % required for inserting images

% for cyrillic and Bulgarian language
\usepackage[utf8]{inputenc}
\usepackage[T2A]{fontenc}
\usepackage[bulgarian]{babel}

% math packages
\usepackage{amsmath}
\usepackage{float}
\usepackage{amsfonts}
\usepackage{amssymb}
\usepackage{amsthm}
\usepackage{cancel}

\usepackage{ragged2e}  % adds FlushLeft

\usepackage{listings}
\usepackage{xcolor}

\definecolor{codegreen}{rgb}{0,0.6,0}
\definecolor{codegray}{rgb}{0.5,0.5,0.5}
\definecolor{codepurple}{rgb}{0.58,0,0.82}
\definecolor{backcolour}{rgb}{0.95,0.95,0.92}

\lstdefinestyle{mystyle}{
    language=Octave,
    backgroundcolor=\color{backcolour},   
    commentstyle=\color{codegreen},
    keywordstyle=\color{magenta},
    numberstyle=\tiny\color{codegray},
    stringstyle=\color{codepurple},
    basicstyle=\ttfamily\footnotesize,
    breakatwhitespace=false,         
    breaklines=true,
    captionpos=b,
    keepspaces=false,                 
    numbersep=5pt,
    showspaces=false,
    showstringspaces=false,
    showtabs=false,
    tabsize=2,
    columns=fullflexible,
}

\lstset{style=mystyle}

\usepackage[colorlinks=true, linkcolor=blue, urlcolor=cyan]{hyperref}

\newtheorem{theorem}{Теорема}
\newtheorem{definition}{Дефиниция}

\title{Линейна алгебра - собствени вектори и собствени стойности}
\author{Ивайло Андреев}
\date{5 май 2026}

\begin{document}
\maketitle

\tableofcontents

\newpage

\begin{FlushLeft}

Нека $V$ е линейно пространство над числово поле $F$ и нека $\dim{V} = n \in \mathbb{N} \setminus \lbrace 0 \rbrace$.

Нека $\varphi \in \text{Hom}{V}$ е линеен оператор.

Тези означения важат за цялата тема.

\section{Определение на характеристичен полином на матрица}

\begin{definition}
(Характеристичен полином и характеристични корени на матрица)

Нека $A \in M_n(F)$.

Характеристичен полином на матрицата $A$ e
\[
P_A(\lambda) = \det(A - \lambda E)
\]

Характеристични корени на матрицата $A$ са корените на характеристичния полином.
\end{definition}

\section{Подобни матрици и
доказателство, че подобните матрици имат едни и същи характеристични полиноми}

\begin{definition}
(Подобна матрица)
Нека $A, B \in M_n(F)$.

Матриците $A$ и $B$ са подобни (и пишем $A \sim B$), ако съществува обратима матрица $T \in M_n(F)$, такава че

$$B = T^{-1}AT$$
\end{definition}

\begin{theorem}
Подобните матрици имат един и същи характеристичен полином.
\end{theorem}

\textbf{Доказателство:}

Нека $A, B \in M_n(F)$ и нека $A \sim B$. Тогава съществува обратима матрица $T$, такава че $B = T^{-1}AT$

Ще покажем, че $P_B(\lambda) = P_A(\lambda)$.

\begin{align*}
P_B(\lambda)
&= \det(B - \lambda E)\\
&= \det(T^{-1}AT - \lambda E)\\
&= \det(T^{-1}AT - \lambda T^{-1} ET)\\
&= \det(T^{-1}(A - \lambda E)T)\\
&= \det(T^{-1})\det(A - \lambda E)\det(T)\\
&= \det(T)^{-1}\det(A - \lambda E)\det(T)\\
&= \det(A - \lambda E)\\
&= P_A(\lambda)\\
\end{align*}

\section{Определение на характеристичен полином на линеен оператор}

\begin{definition}
(Характеристичен полином и характеристични корени на линеен оператор)

Характеристичен полином на линейния оператор $\varphi$ е характеристичния полином на матрицата на $\varphi$ в който и да е базис на $V$.

Характеристични корени на линейния оператор $\varphi$ са корените на характеристичния полином.
\end{definition}

\section{Определение на
собствен вектор и собствена стойност на линеен оператор}

\begin{definition}
(Собствен вектор и собствена стойност)

Нека $\vec{v}\in V$.

Наричаме вектора $\vec{v}$ собствен вектор, съответстващ на собствената стойност $\lambda \in F$, ако

\begin{enumerate}
    \item $\vec{v} \ne \vec{0}$
    \item $\varphi(\vec{v}) = \lambda \vec{v}$
\end{enumerate}
\end{definition}

\section{Необходимо и достатъчно
условие характеристичен корен на линеен оператор да е собствена стойност в
крайномерно линейно пространство}

\begin{theorem}
Характеристични корени на линейния оператор $\varphi$, които са в полето $F$, са (единствените) собствени стойности на линейния оператор $\varphi$.
\end{theorem}

\textbf{Доказателство:}

Нека $\{e_i\}_{i=1}^n$ е фиксиран базис на $V$.

Нека $A$ е матрицата на $\varphi$ в базиса $\{e_i\}_{i=1}^n$.

\textbf{Права посока}

Нека $\lambda \in F$ е собствена стойност, съответстваща на ненулевия собствен вектор $\vec{v} \neq \vec{0}$. Ще покажем, че $\lambda$ е корен на характеристичния полином на $\varphi$.

По дефиниция

$$\varphi(\vec{v}) = \lambda \vec{v}$$

Нека координатите на $\vec{v}$ в базиса $\{e_i\}_{i=1}^n$ са $\vec{\alpha}$

$$\vec{v} = \sum_{i=1}^n \alpha_i e_i$$

Знаем, че $\vec{v}$ е собствен вектор, в частност ненулев, и съответно координатите му не са нулев вектор. Тоест

$$\vec{\alpha} \ne \vec{0}$$

Нека координатите на $\varphi(\vec{v})$ в базиса $\{e_i\}_{i=1}^n$ са $\vec{\beta}$

$$\varphi(\vec{v}) = \sum_{i=1}^n \beta_i e_i$$

Знаем, че $\vec{v}$ е собствен вектор и е изпълнено равенството

$$\varphi(\vec{v}) = \lambda \vec{v}$$

$$\sum_{i=1}^n \beta_i e_i = \lambda \sum_{i=1}^n \alpha_i e_i$$

$$\sum_{i=1}^n \beta_i e_i = \sum_{i=1}^n (\lambda\alpha_i) e_i$$

Приравнявайки подобните коефициенти пред базисните вектори получаваме

$$\vec{\beta} = \lambda \vec{\alpha}$$

От друга страна от теоремата за линейни оператори знаем, че

$$\vec{\beta} = A \vec{\alpha}$$

Приравняваме и получаваме

$$A \vec{\alpha} = \lambda \vec{\alpha}$$

$$A \vec{\alpha} - \lambda \vec{\alpha} = 0$$

$$A \vec{\alpha} - \lambda E \vec{\alpha} = 0$$

$$(A  - \lambda E) \vec{\alpha} = 0$$

Получаваме хомогенна система линейни уравнения. Вече споменахме, че $\vec{\alpha} \ne \vec{0}$ и съответно системата притежава нетривиални решения. Тогава детерминантата на матрицата $A - \lambda E$ е нула.

$$\det{(A - \lambda E)} = 0$$

$$P_A(\lambda) = 0$$

Така получаваме, че $\lambda$, която по начало казахме, че е собствена стойност на линейния оператор, също е и корен на характеристичния полином на линейния оператор.

\textbf{Обратна посока}

Нека $\lambda \in F$ е корен на характеристичния полином на линейния оператор $\varphi$. Ще покажем, че $\lambda$ е собствена стойност на $\varphi$.

По дефиниция

$$\det(A - \lambda E) = 0$$

Разглеждаме хомогенната система линейни уравнения спрямо неизвестния вектор-стълб $\vec{x}$

$$(A - \lambda E)\vec{x} = \vec{0}$$

Тъй като детерминантата на матрицата на системата е нула, системата притежава ненулево решение $\vec{\alpha} \neq \vec{0}$, такова че

$$A \vec{\alpha} = \lambda \vec{\alpha}$$

Нека $\vec{v} = \sum_{i=1}^n \alpha_i e_i$ е векторът от $V$, чийто координати в базиса $\{e_i\}_{i=1}^n$ са точно елементите на вектора $\vec{\alpha}$. Знаем, че $\vec{\alpha} \neq \vec{0}$, то и $\vec{v} \neq \vec{0}$.

От матричното представяне на линейния оператор знаем, че ако координатите на $\varphi(\vec{v})$ в базиса $\{e_i\}_{i=1}^n$ са $\vec{\beta}$, то е изпълнено равенството 

$$\vec{\beta} = A \vec{\alpha}$$

След приравняване

$$\vec{\beta} = A \vec{\alpha} = \lambda \vec{\alpha}$$

Това означава, че векторът $\varphi(\vec{v})$ има координати $\lambda \vec{\alpha}$, т.е.

$$\varphi(\vec{v}) = \sum_{i=1}^n (\lambda \alpha_i) e_i = \lambda \sum_{i=1}^n \alpha_i e_i = \lambda \vec{v}$$

Тъй като $\vec{v} \neq \vec{0}$ и $\varphi(\vec{v}) = \lambda \vec{v}$, то по дефиниция $\lambda$ е собствена стойност на линейния оператор $\varphi$.

С това теоремата е доказана в двете посоки.

\section{Доказателство, че собствените вектори на линеен
оператор, съответстващи на различни собствени стойности са линейно независими}

\begin{theorem}\label{theorem_linear_independency}
Собствените вектори, съответстващи на различни собствени стойности, са линейно независими.
\end{theorem}

\textbf{Доказателство:}

Нека $\lambda_1, \dots, \lambda_k$ са различни собствени стойности, съответстващи на собствените вектори $\vec{v}_1, \dots, \vec{v}_k$.

Индукция по $k$.

\textbf{База}: $k= 1$

По определение собствения вектор е ненулев и знаем, че система от един ненулев вектор е линейно независима.

\textbf{Индукционно предположение}: $k = n$

Нека $\lambda_1, \dots, \lambda_n$ са различни собствени стойности, съответстващи на линейно независимите собствените вектори $\vec{v}_1, \dots, \vec{v}_n$. Допускаме, че твърдението е вярно за всяка система от $n$ вектора.

\textbf{Стъпка}: $k = n + 1$

Разглеждаме равенството

\begin{equation}
    \beta_1 \vec{v}_1 + \dots + \beta_n \vec{v}_n + \beta_{n+1} \vec{v}_{n+1} = \vec{0} \label{bofore_phi}
\end{equation}

Ще покажем, че всички коефициенти $\beta_i = 0 \quad i \in \{1 \dots n+1\}$ от линейната комбинация са нули, откъдето ще следва, че векторите са линейно независими.

$$\varphi(\beta_1 \vec{v}_1 + \dots + \beta_n \vec{v}_n + \beta_{n+1} \vec{v}_{n+1}) = \varphi(\vec{0})$$

$$\varphi(\beta_1 \vec{v}_1) + \dots + \varphi(\beta_n \vec{v}_n) + \varphi(\beta_{n+1} \vec{v}_{n+1}) = \varphi(\vec{0})$$

\begin{equation}
    \lambda_1\beta_1 \vec{v}_1 + \dots + \lambda_n\beta_n \vec{v}_n + \lambda_{n+1}\beta_{n+1} \vec{v}_{n+1} = \vec{0} \label{after_phi}
\end{equation}

Имаме уравненията \eqref{bofore_phi} и \eqref{after_phi}. Целим се да елиминираме терма $\vec{v}_{n+1}$, за да можем да стигнем до уравнение, където да можем да използваме индукционното предположение. Ще постигнем това като умножим равенството \eqref{bofore_phi} с $\lambda_{n+1}$ и го извадим от равенството \eqref{after_phi}.

$$(\lambda_1 - \lambda_{n+1})\beta_1 \vec{v}_1 + \dots + (\lambda_n - \lambda_{n+1})\beta_n \vec{v}_n + (\lambda_{n+1} - \lambda_{n+1})\beta_{n+1} \vec{v}_{n+1} = \vec{0}$$

\begin{equation}
    (\lambda_1 - \lambda_{n+1})\beta_1 \vec{v}_1 + \dots + (\lambda_n - \lambda_{n+1})\beta_n \vec{v}_n = \vec{0} \label{after_subtraction}
\end{equation}

От индукционното предположение знаем, че $\vec{v}_1, \dots, \vec{v}_n$ са линейно независими. Това означава, че коефициентите в линейната комбинация, даваща нула, също са нули. В \eqref{after_subtraction} коефициентите $\lambda_i - \lambda_{n+1} \ne 0 \quad i \in \{1 \dots n\}$ са ненулеви, защото собствените стойности са различни по условие. Тогава остава $\beta_i = 0 \quad i \in \{1 \dots n\}$.

Заместваме в \eqref{bofore_phi} с нулевите коефициенти

$$\beta_{n+1} \vec{v}_{n+1} = \vec{0}$$

Векторът $\vec{v}_{n+1}$ е собствен и в частност ненулев, откъдето $\beta_{n+1} = 0$.

Показахме, че $\beta_i = 0 \quad i \in \{1 \dots n+1\}$, с което теоремата е доказана.

\section{Линеен оператор с прост спектър (с доказателство)}

\begin{definition}
Линеен оператор има прост спектър, ако има $n$ различни собствени стойности.
\end{definition}

\begin{theorem}
Съществува базис на линеен оператор с прост спектър $\varphi \in \text{Hom}{V}$, в който матрицата на $\varphi$ да е диагонална.
\end{theorem}

\textbf{Доказателство:}

Линейният оператор $\varphi$ е с прост спектър и съответно притежава $n$ на брой различни собствени стойности. Нека това са $\lambda_1$, $\dots$, $\lambda_n$. Нека на тези собствени стойности съответстват собствените вектори $\vec{v}_1$, $\dots$, $\vec{v}_n$. Според теорема \ref{theorem_linear_independency} $\vec{v}_1$, $\dots$, $\vec{v}_n$ са линейно независими и тъй като са $n$ на брой образуват базис на $V$.

Знаем, че $\varphi(\vec{v}_i) = \lambda_i \vec{v}_i$, понеже векторите са собствени. От определението за матрица на линеен оператор следва, че матрицата на $\varphi$ в базиса на собствените вектори има вида

$$
D =
\begin{pmatrix}
    \lambda_1 & & 0\\
    & \ddots &\\
    0 & & \lambda_n
\end{pmatrix}
$$

което е диагонална матрица, с което теоремата е доказана.

\section{Определение за ф-инвариантно
подпространство на линейно пространство}

\begin{definition}
Нека $U$ е подпространство на $V$.

Подпространството $U$ e $\varphi$-инвариантно, ако за всеки вектор $\vec{u} \in U$ е изпълнено, че $\varphi(\vec{u}) \in U$
\end{definition}

\section{Основни примери за ф-инвариантни
подпространства (с доказателство)}

\subsection{Пример 1}

Очевидно $U = V$ е $\varphi$-инвариантно подпространство.

$$\dots$$

Подпространството $U = \lbrace \vec{0} \rbrace$ е $\varphi$-инвариантно за всеки линеен оператор $\varphi$, защото линейните оператори имат свойството, че $\varphi(\vec{0}) = \vec{0}$.

\subsection{Пример 2}

Подпространството $\text{Ker}\varphi$ е $\varphi$-инвариантно подпространство.

Нека $\vec{u} \in \text{Ker}\varphi$. Тогава $\varphi(\vec{u}) = \vec{0}$

$$\varphi(\varphi(\vec{u})) = \varphi(\vec{0}) = \vec{0} \in \text{Ker}\varphi$$

$$\dots$$

Подпространството $\text{Im}\varphi$ е $\varphi$-инвариантно подпространство.

Нека $\vec{b} \in \text{Im}\varphi$. Тогава съществува $\vec{a}$, за който $\varphi(\vec{a}) = \vec{b}$.

$$\varphi(\varphi(\vec{a})) = \varphi(\vec{b}) \in \text{Im}\varphi$$

\subsection{Пример 3}

Подпространството $U = \{\vec{v} \in V \space | \space \varphi(\vec{v}) = \lambda \vec{v}\}$, където $\lambda$ е фиксирана собствена стойност, е $\varphi$-инвариантно подпространство.

Нека $\vec{u} \in U$. Тогава $\varphi(\vec{u}) = \lambda \vec{u}$.

$$\varphi(\varphi(\vec{u})) = \varphi(\lambda \vec{u}) = \lambda\varphi(\vec{u})$$

$\vec{w} = \varphi(\vec{u}) \in U$ и също така $\lambda \vec{w} \in U$

$$\dots$$

Подпространството $U = l(\vec{v})$, където $\vec{v}$ е фиксиран собствен вектор, е $\varphi$-инвариантно подпространство.

Нека $\vec{u} \in U$. Тогава съществува скалар $\alpha \in F$, такъв че $\vec{u} = \alpha \vec{v}$.

Знаем, че $\vec{v}$ е собствен вектор и съответно $\varphi(\vec{v}) = \lambda \vec{v}$.

$$\varphi(\vec{u}) = \varphi(\alpha \vec{v}) = \alpha\varphi(\vec{v}) = \alpha(\lambda\vec{v}) = (\alpha\lambda)\vec{v} \in U$$

\section{Теореми за съществуване на ф-инвариантни
подпространства на линейно пространство над полето на реалните и на комплексните числа (без доказателство)}

\begin{theorem}
Нека $V$ е линейно пространство над $\mathbb{R}$ и нека $\dim{V} = n \in \mathbb{N} \setminus \lbrace 0 \rbrace$ и нека $\varphi \in \text{Hom}{V}$. Тогава $\varphi$ има едномерно или двумерно $\varphi$-инвариантно подпространство.
\end{theorem}

\begin{theorem}
Нека $V$ е линейно пространство над $\mathbb{C}$ и нека $\dim{V} = n \in \mathbb{N} \setminus \lbrace 0 \rbrace$ и нека $\varphi \in \text{Hom}{V}$. Тогава за всяка собствена стойност $\lambda \in \mathbb{C}$ съществува собствен вектор $\vec{v}$, такъв че едномерното подпространство на линейната му обвивка $l(v)$ е $\varphi$-инвариантно.
\end{theorem}

\end{FlushLeft}

\end{document}